\subsection{闭子集的余维数}

定义 3. --- 设 X 是一个拓扑空间。

\begin{enumerate}
\def\labelenumi{\alph{enumi})}
\item
  如果 Y 是 X 的闭不可约子集，则称 Y 在 X 中的余维数为所有以 Y 为最小元的 X 的不可约闭子集链的长度在 \(\overline{\mathbf{R}}\) 中的上确界。
\item
  如果 Y 是 X 的闭子集，则称 Y 在 X 中的余维数，并记为 \(\text{codim}(Y, X)\)，为 Y 的不可约分支在 X 中的余维数在 \(\overline{\mathbf{R}}\) 中的下确界。
\end{enumerate}

注。--- 1) 因此，X 的闭子集 Y 的余维数是 Y 的不可约闭子集的余维数的下确界。我们有 \(\mathrm{codim}(\emptyset, \mathrm{X}) = +\infty\)；如果 X 非空，则 \(\mathrm{codim}(\mathrm{X}, \mathrm{X}) = 0\)。X 的每个非空闭子集都含有一个不可约闭子集（II, § 4, no. 1, prop. 5）；因此，闭子集 Y 在 X 中的余维数总是一个非负整数或 \(+\infty\)。当且仅当 Y 含有 X 的一个不可约分支时，它为零。

\begin{enumerate}
\def\labelenumi{\arabic{enumi})}
\setcounter{enumi}{1}
\tightlist
\item
  如果 Y 是 X 的非空闭子集，则 codim(Y, X) ≤ dim(X)。此外，dim(X) = sup codim(Y, X)，其中 Y 遍历不可约闭子集的集合。
\end{enumerate}

如果 Y 和 Y' 是 X 的两个闭子集且 Y' ⊂ Y，则 codim(Y, X) ≤ codim(Y', X)。

\begin{enumerate}
\def\labelenumi{\arabic{enumi})}
\setcounter{enumi}{2}
\tightlist
\item
  设 Y 是拓扑空间 X 的一个闭子集，且 \((\mathrm{X}_{\alpha})_{\alpha\in\mathrm{A}}\)（分别为 \((\mathrm{Y}_{\beta})_{\beta\in\mathrm{B}}\)）是 X（分别为 Y）的不可约分支族。对每个 \(\beta\in B\)，记 \(\mathrm{A}(\beta)\) 为由满足 \(\alpha\in A\) 且 \(Y_{\beta}\subset X_{\alpha}\) 的指标组成的集合。由于 X 的每个不可约子集都包含于某个 \(X_{\alpha}\) 中（II, § 4, no. 1, prop. 5），由定义 3 得：
\end{enumerate}

\[
\operatorname{codim} (Y, X) = \inf _ {\beta \in B} \sup _ {\alpha \in A (\beta)} \operatorname{codim} (Y _ {\beta}, X _ {\alpha}).
\]

\begin{enumerate}
\def\labelenumi{\arabic{enumi})}
\setcounter{enumi}{3}
\tightlist
\item
  设 \((\mathbf{Y}_i)_{i\in I}\) 是 X 的闭子集的一个有限族，且 \(\mathbf{Y} = \bigcup_{i\in I}\mathbf{Y}_i\)；有
\end{enumerate}

\[
\operatorname{codim} (\mathrm{Y}, \mathrm{X}) = \inf _ {i \in \mathrm{I}} \operatorname{codim} (\mathrm{Y} _ {i}, \mathrm{X}).
\]

事实上，Y 的每个不可约分支都包含于某个 \(Y_{i}\) 中。

命题 3. --- 设 X 是一个拓扑空间。

\begin{enumerate}
\def\labelenumi{\alph{enumi})}
\tightlist
\item
  对 X 的每个非空闭子集 Y，都有
\end{enumerate}

\[
\dim (\mathrm{Y}) + \operatorname{codim} (\mathrm{Y}, \mathrm{X}) \leqslant \dim (\mathrm{X}).
\]

\begin{enumerate}
\def\labelenumi{\alph{enumi})}
\setcounter{enumi}{1}
\tightlist
\item
  如果 Y、Z、T 是 X 的闭子集且 Y ⊂ Z ⊂ T，则
\end{enumerate}

\[
\operatorname{codim} (\mathrm{Y}, \mathrm{Z}) + \operatorname{codim} (\mathrm{Z}, \mathrm{T}) \leqslant \operatorname{codim} (\mathrm{Y}, \mathrm{T}).
\]

只需证明断言 \(a)\)，在 \(\dim(X)\) 有限的情形下。此时，\(\dim(Y)\) 和 \(\operatorname{codim}(Y, X)\) 均有限。存在一条 \(Y_0 \subset \ldots \subset Y_n\) 的 \(Y\) 的不可约闭子集链，长度为 \(n = \dim(Y)\)，以及一条 \(Y_n \subset \ldots \subset Y_{n+p}\) 的 \(X\) 的不可约闭子集链，长度为 \(p \geqslant \operatorname{codim}(Y, X)\)。由此 \(\dim(X) \geqslant n + p\)，从而得到结论。为证明 \(a)\)\(b)\)，可假设 \(Y\) 不可约。由于有 \(\operatorname{codim}(Y, Z) \leqslant \operatorname{codim}(Y, T)\)，当 \(\operatorname{codim}(Y, Z) = +\infty\) 时不等式成立。否则，设 \(Z_0\) 是 \(Z\) 的一个不可约分支，包含 \(Y\)，且满足 \(\operatorname{codim}(Y, Z) = \operatorname{codim}(Y, Z_0)\)。有 \(\operatorname{codim}(Z, T) \leqslant \operatorname{codim}(Z_0, T)\)，并且如上可见，\(\operatorname{codim}(Y, Z_0) + \operatorname{codim}(Z_0, T) \leqslant \operatorname{codim}(Y, T)\)，从而 \(b)\)。

定义 4. --- 如果对 X 的每一对满足 Y ⊂ Z 的不可约闭子集 (Y, Z)，以 Y 和 Z 为端点的每条不可约闭子集饱和链的长度都等于 codim(Y, Z)，则称拓扑空间 X 为串联的。

等价地说，对于 X 的每一对满足 codim(Y, Z) 有限的不可约闭子集 (Y, Z)，以 Y 和 Z 为端点的所有饱和链具有相同长度；而对于满足 codim(Y, Z) = +∞ 的每一对 (Y, Z)，不存在以 Y 和 Z 为端点的饱和链。

串联空间的每个闭子空间都是串联的。一个空间为串联的，当且仅当它的不可约分支都是串联的。

命题 4. --- 设 X 是一个拓扑空间。X 为串联的，当且仅当对 X 的每个满足 Y ⊂ Z ⊂ T 的不可约闭子集三元组 (Y, Z, T)，都有：

\[
\operatorname{codim} (\mathrm{Y}, \mathrm{T}) = \operatorname{codim} (\mathrm{Y}, \mathrm{Z}) + \operatorname{codim} (\mathrm{Z}, \mathrm{T}).
\]

设 X 是串联的。根据命题 3 b)，只需在 codim(Y, Z) 和 codim(Z, T) 均有限时证明该关系。将一条以 Y 和 Z 为端点、长度为 codim(Y, Z) 的不可约闭子集饱和链，与一条以 Z 和 T 为端点、长度为 codim(Z, T) 的不可约闭子集饱和链首尾相接，便得到一条以 Y 和 T 为端点、长度为 codim(Y, Z) + codim(Z, T) 的饱和链。但是，由于 X 是串联的，这个长度必然等于 codim(Y, T)。

反之，假设有 \(\operatorname{codim}(\mathbf{Y},\mathbf{T})=\operatorname{codim}(\mathbf{Y},\mathbf{Z})+\operatorname{codim}(\mathbf{Z},\mathbf{T})\)，并且该等式对满足 \(\mathbf{Y}\subset\mathbf{Z}\subset\mathbf{T}\) 的所有不可约闭子集均成立；我们证明 X 是串联的。为此，我们对整数 \(n\geqslant0\) 作归纳，证明对于 X 的每条不可约闭子集饱和链 \(Z_{0}\subset\ldots\subset Z_{n}\)，都有 \(\operatorname{codim}(Z_{0},Z_{n})=n\)。当 \(n=0\) 时显然成立。设 \(n>0\)，并假设长度 \(\leqslant n-1\) 的链都满足该性质。如果 \(Z_{0}\subset\ldots\subset Z_{n}\) 是长度为 \(n\) 的饱和链，则 \(Z_{0}\subset\ldots\subset Z_{n-1}\) 是长度为 \(n-1\) 的饱和链，因而 \(\operatorname{codim}(Z_{0},Z_{n-1})=n-1\)。根据关于 X 的假设，有 \(\operatorname{codim}(Z_{0},Z_{n})=\operatorname{codim}(Z_{0},Z_{n-1})+\operatorname{codim}(Z_{n-1},Z_{n})=(n-1)+1=n\)。

推论。— 设 X 是一个有限维不可约拓扑空间。X 为串联的，当且仅当对 X 的每一对满足 Y ⊂ Z 的不可约闭子集 (Y, Z)，有 codim(Y, X) = codim(Y, Z) + codim(Z, X)。

由命题 4，必要性成立。反之，假设该条件成立，并记 \(c(Z)\) 为整数 \(\operatorname{codim}(Z,X)\)，其中 \(Z\) 遍历 \(X\) 的每个不可约闭子集。如果 \(Y,Z,T\) 是 \(X\) 的三个不可约闭子集，且 \(Y\subset Z\subset T\)，则有

\[
\begin{array}{r l} \operatorname{codim} (\mathrm{Y}, \mathrm{Z}) + \operatorname{codim} (\mathrm{Z}, \mathrm{T}) & = (c (\mathrm{Y}) - c (\mathrm{Z})) + (c (\mathrm{Z}) - c (\mathrm{T})) \\ & = c (\mathrm{Y}) - c (\mathrm{T}) \\ & = \operatorname{codim} (\mathrm{Y}, \mathrm{T}), \end{array}
\]

且由命题 4，X 是串联的。

命题 5. — 设 X 是一个有限维拓扑空间。假设 X 的所有不可约闭子集极大链具有相同长度。则 X 是串联的；对 X 的每个不可约闭子集 Z，有

\[
\operatorname{codim} (Z, X) = \dim (X) - \dim (Z);
\]

对 X 的每一对满足 Y ⊂ Z 的不可约闭子集 (Y, Z)，有

\[
\operatorname{codim} (\mathrm{Y}, \mathrm{Z}) = \dim (\mathrm{Z}) - \dim (\mathrm{Y}).
\]

设 Y 和 Z 是 X 的两个不可约闭子集且 \(Y \subset Z\)。设 \(Y_{0} \subset \ldots \subset Y_{p}\) 是一条满足 \(Y_{p} = Y\) 且 \(p = \dim(Y)\) 的链，\(Z_{0} \subset \ldots \subset Z_{q}\) 是一条满足 \(Z_{0} = Z\) 且 \(q = \text{codim}(Z, X)\) 的链。对于每条满足 \(T_{0} \subset \ldots \subset T_{r}\)、\(T_{0} = Y\) 且 \(T_{r} = Z\) 的饱和链，该链

\[
\mathrm{Y} _ {0} \subset \dots \subset \mathrm{Y} _ {p - 1} \subset \mathrm{T} _ {0} \subset \dots \subset \mathrm{T} _ {r} \subset \mathrm{Z} _ {1} \dots \subset \mathrm{Z} _ {q}
\]

它是极大的，且长度为 \(p + q + r\)；根据关于 X 的假设，有 \(p + q + r = \dim(X)\)。令 \(r = \dim(X) - \dim(Y) - \text{codim}(Z, X)\)。由此 X 是串联的，并且对于上述 Y、Z，有

\[
\dim (\mathrm{Y}) + \operatorname{codim} (\mathrm{Y}, \mathrm{Z}) = \dim (\mathrm{X}) - \operatorname{codim} (\mathrm{Z}, \mathrm{X}).
\]

取 Y = Z，可见第二项等于 dim(Z)，从而得到命题。

